\section{Prepared Statement dan Parameterized Query: Mencegah SQL Injection}

\textbf{SQL injection} adalah serangan paling berbahaya pada aplikasi web yang menggunakan database. Penyerang menyisipkan kode SQL melalui input pengguna, memanipulasi query untuk mengakses, mengubah, atau menghapus data yang tidak seharusnya. \textbf{Prepared statement} adalah pertahanan utama: query dikirim ke database sebagai template dengan placeholder (\code{?}), lalu data dikirim terpisah. Database memperlakukan data sebagai literal, bukan perintah SQL \cite{phpManual}.

\begin{contoh}
\textbf{Studi Kasus: SQL Injection pada Form Login}

\textbf{Query rentan}: \code{"SELECT * FROM users WHERE user='\$u' AND pass='\$p'"}.

Jika penyerang memasukkan \code{admin' OR '1'='1} sebagai username, query menjadi: \code{SELECT * FROM users WHERE user='admin' OR '1'='1' AND pass='...'} --- kondisi \code{'1'='1'} selalu benar, sehingga autentikasi terlewati.

\textbf{Solusi}: prepared statement memisahkan template dari data, sehingga input berbahaya diperlakukan sebagai string biasa, bukan perintah SQL.
\end{contoh}

\begin{webcode}[language=PHP, caption={Prepared statement untuk login aman}]
<?php
$email    = $_POST['email'] ?? '';
$password = $_POST['password'] ?? '';

// AMAN: prepared statement
$stmt = $conn->prepare("SELECT id, nama, password FROM users WHERE email = ?");
$stmt->bind_param("s", $email);
$stmt->execute();
$result = $stmt->get_result();
$user   = $result->fetch_assoc();

if ($user && password_verify($password, $user['password'])) {
  session_regenerate_id(true);
  $_SESSION['user_id'] = $user['id'];
  $_SESSION['nama']    = $user['nama'];
  header("Location: dashboard.php");
} else {
  echo "Email atau password salah.";
}
$stmt->close();
?>
\end{webcode}

